\section*{第 \ref{ch:b1:digestion-absorption} 章——消化与吸收}

\begin{solution}{exo:b1:digestion-absorption:1}
口腔：咀嚼，唾液\omterm{def:b1:carbohydrates:polysaccharide}{淀粉}酶。食管：运送。胃：
胃酸、胃蛋白酶，储存与混合。小肠：胆汁与
胰酶完成\omterm{def:b1:digestion-absorption:nutrition}{消化}；\omterm{def:b1:digestion-absorption:nutrition}{吸收}。大肠：回收水分和盐分，
微生物\omterm{def:b1:respiration-fermentation:fermentation}{发酵}。直肠：
储存与排遗。
\end{solution}

\begin{solution}{exo:b1:digestion-absorption:2}
胃中的胃蛋白酶在 \omterm{def:b1:water-small-molecules:ph}{pH} 2 下起作用；胰腺分泌的胰蛋白酶、糜蛋白酶和
羧肽酶在小肠中、\omterm{def:b1:water-small-molecules:ph}{pH} 8 下起作用；刷状缘肽酶在 \omterm{def:b1:water-small-molecules:ph}{pH} 7 下起作用。
它们以酶原形式分泌，以免\omterm{def:b1:digestion-absorption:nutrition}{消化}产生它们的\omterm{def:b1:cell-unit-of-life:cell}{细胞}；
仅在肠腔中才被激活。
\end{solution}

\begin{solution}{exo:b1:digestion-absorption:3}
葡萄糖：钠同向转运入\omterm{def:b1:cell-unit-of-life:cell}{细胞}，再由\omterm{def:b1:membranes-transport:transporters}{载体蛋白}运出，
经门静脉至肝脏。\omterm{def:b1:proteins:aminoacid}{氨基酸}：钠同向转运入\omterm{def:b1:cell-unit-of-life:cell}{细胞}，\omterm{def:b1:membranes-transport:transporters}{载体蛋白}运出，经门静脉。\omterm{def:b1:lipids:fattyacid}{脂肪
酸}：从胶束中扩散入\omterm{def:b1:cell-unit-of-life:cell}{细胞}，再酯化，包装成乳糜微粒，
胞吐进入乳糜管和\omterm{def:b1:mammal-organization:compartments}{淋巴液}。
\end{solution}

\begin{solution}{exo:b1:digestion-absorption:4}
胆盐将\omterm{def:b1:lipids:triglyceride}{脂肪}乳化为微米级的液滴，并以胶束形式将
\omterm{def:b1:lipids:triglyceride}{脂肪}酶的产物运送至\omterm{def:b1:digestion-absorption:surface}{肠上皮细胞}。若无胆盐，\omterm{def:b1:lipids:triglyceride}{脂肪}
会以大液滴存在，几乎没有供\omterm{def:b1:lipids:triglyceride}{脂肪}酶作用的表面，
释放出的\omterm{def:b1:lipids:fattyacid}{脂肪酸}无法输送到膜上，\omterm{def:b1:lipids:triglyceride}{脂肪}便未经\omterm{def:b1:digestion-absorption:nutrition}{消化}
直接通过。
\end{solution}

\begin{solution}{exo:b1:digestion-absorption:5}
$100/162 = \qty{0.62}{mol}$ 的葡萄糖单位，因此约
有 \qty{0.62}{mol} 的键和 \qty{0.62}{mol} 的水（\qty{11}{g}）；
$620/180 = \qty{3.4}{mmol/min}$，即葡萄糖的\omterm{def:b1:digestion-absorption:nutrition}{吸收}速率。
\end{solution}

\begin{solution}{exo:b1:digestion-absorption:6}
胃酸促使释放促胰液素，后者使胰腺分泌碳酸氢盐，
用以中和食糜；\omterm{def:b1:lipids:triglyceride}{脂肪}和肽促使释放缩胆囊素，它使胆囊收缩
（胆汁流入肠道），刺激胰酶分泌，并抑制胃排空，
使十二指肠接收食糜的速度
仅与其中和并\omterm{def:b1:digestion-absorption:nutrition}{消化}食糜的速度相当，
即不超过其处理能力。
\end{solution}

\begin{solution}{exo:b1:digestion-absorption:7}
葡萄糖被逆梯度浓缩至八倍，属于\omterm{def:b1:membranes-transport:transporters}{主动运输}。
若无钠，运输在达到平衡时即停止：摄取
与钠梯度偶联（钠–葡萄糖\omterm{prop:b1:membranes-transport:secondary}{同向转运体}），
而并非直接与 \omterm{def:b1:water-small-molecules:atp}{ATP} 偶联。
\end{solution}

\begin{solution}{exo:b1:digestion-absorption:8}
乳糜微粒直径约一微米，不能穿过紧密的毛细血管壁，
但能进入开口的乳糜管；\omterm{def:b1:mammal-organization:compartments}{淋巴液}在颈部汇入血液，
因此\omterm{def:b1:lipids:triglyceride}{脂肪}先到达全身，之后才到肝脏，
否则肝脏会先将其全部摄取。若乳糜微粒直接进入门静脉血，
它们会堵塞肝血窦，并使肝脏在每餐后处理\omterm{def:b1:lipids:triglyceride}{脂肪}的
容量饱和。
\end{solution}

\begin{solution}{exo:b1:digestion-absorption:9}
唾液使纤维湿润并漂浮，以便反刍；其中的碳酸氢盐和
磷酸盐作为\omterm{def:b1:water-small-molecules:buffer}{缓冲液}，缓冲微生物产生的酸，使瘤胃
保持在 \omterm{def:b1:water-small-molecules:ph}{pH} 6.5 左右；并且它从血液中带来尿素，微生物将其
转化为氨，再合成自身的\omterm{def:b1:proteins:peptide}{蛋白质}，使氮得以回收利用，
否则这些氮将随尿液流失。
\end{solution}

\begin{solution}{exo:b1:digestion-absorption:10}
若缺乏胰\omterm{def:b1:carbohydrates:polysaccharide}{淀粉}酶、蛋白酶、\omterm{def:b1:lipids:triglyceride}{脂肪}酶和碳酸氢盐：\omterm{def:b1:carbohydrates:polysaccharide}{淀粉}
仅部分由唾液和刷状缘\omterm{def:b1:digestion-absorption:nutrition}{消化}；\omterm{def:b1:proteins:peptide}{蛋白质}
经胃蛋白酶切割，但不能达到可\omterm{def:b1:digestion-absorption:nutrition}{吸收}的大小；\omterm{def:b1:lipids:triglyceride}{脂肪}完全不被
\omterm{def:b1:digestion-absorption:nutrition}{消化}，以体积大、色浅、油腻的粪便排出。\omterm{def:b1:carbohydrates:monosaccharide}{糖类}和\omterm{def:b1:carbohydrates:glycosidic}{二糖}
（刷状缘酶仍存在）、部分\omterm{def:b1:carbohydrates:polysaccharide}{淀粉}及小肽仍可被处理；
而\omterm{def:b1:lipids:triglyceride}{脂肪}和大部分\omterm{def:b1:proteins:peptide}{蛋白质}则不能。
\end{solution}

\begin{solution}{exo:b1:digestion-absorption:11}
盲肠粪便中含有盲肠产生的微生物\omterm{def:b1:proteins:peptide}{蛋白质}
和 B 族维生素。盲肠位于小肠之后，
因此在第一次通过时无法被\omterm{def:b1:digestion-absorption:nutrition}{消化}；
若再次食入，它们会经过胃和小肠并被\omterm{def:b1:digestion-absorption:nutrition}{吸收}。
牛在胃之前\omterm{def:b1:respiration-fermentation:fermentation}{发酵}，其微生物在下游被\omterm{def:b1:digestion-absorption:nutrition}{消化}，
无需二次通过。马的盲肠同样位于下游，
而马既不食用自己的粪便，也没有其他途径，
其微生物\omterm{def:b1:proteins:peptide}{蛋白质}因此损失。
\end{solution}

\begin{solution}{exo:b1:digestion-absorption:12}
微生物以\omterm{def:b1:lipids:fattyacid}{脂肪酸}的形式为牛提供\omterm{def:b1:carbohydrates:polysaccharide}{纤维素}中的能量，
并由贫劣牧草甚至尿素合成\omterm{def:b1:proteins:peptide}{蛋白质}，还提供所有 B 族维生素；
牛则为它们提供温暖、有缓冲作用、厌氧的\omterm{def:b1:respiration-fermentation:fermentation}{发酵}瓮，
保持充满并不断搅动，并每天分泌一百升唾液。这种设计的代价是：
十分之一的能量以甲烷形式逸出，无葡萄糖被\omterm{def:b1:digestion-absorption:nutrition}{吸收}，
故肝脏必须全部由丙酸盐生成葡萄糖，\omterm{def:b1:digestion-absorption:nutrition}{消化}需要三天，
而\omterm{def:b1:respiration-fermentation:fermentation}{发酵}瓮及其内容物的重量占动物体重的五分之一——
牛不能冲刺，狮子却能。
\end{solution}

\begin{solution}{pb:b1:digestion-absorption:1}
\textbf{1.} 纤维 \qty{5}{kg}，\omterm{def:b1:carbohydrates:polysaccharide}{淀粉}和\omterm{def:b1:carbohydrates:monosaccharide}{糖类} \qty{2}{kg}，\omterm{def:b1:proteins:peptide}{蛋白质}
\qty{1.5}{kg}。
\textbf{2.} $0.7\times 5 + 2 = \qty{5.5}{kg}$。
\textbf{3.} $5500\times 17.5 = \qty{96}{MJ}$；挥发性\omterm{def:b1:lipids:fattyacid}{脂肪酸} $0.75\times 96 =
\qty{72}{MJ}$。
\textbf{4.} $0.08\times 96 = \qty{7.7}{MJ}$；$7700/55 = \qty{140}{g}$
甲烷，约 \qty{200}{L}。
\textbf{5.} $0.15\times 5500 = \qty{825}{g}$。
\textbf{6.} 大部分膳食\omterm{def:b1:proteins:peptide}{蛋白质}被瘤胃中的微生物降解，
并重新合成为微生物\omterm{def:b1:proteins:peptide}{蛋白质}；微生物随后流入
皱胃和小肠，在那里与任何肉类一样被\omterm{def:b1:digestion-absorption:nutrition}{消化}。
\textbf{7.} 纤维的 \qty{30}{\%}，即 \qty{1.5}{kg}：部分在结肠中\omterm{def:b1:respiration-fermentation:fermentation}{发酵}，
大部分被排出。
\textbf{8.} 逸过瘤胃的膳食\omterm{def:b1:proteins:peptide}{蛋白质} $0.4\times 1500 = \qty{600}{g}$，
可\omterm{def:b1:digestion-absorption:nutrition}{消化}部分 \qty{360}{g}；微生物\omterm{def:b1:proteins:peptide}{蛋白质} \qty{825}{g}；合计 \qty{1185}{g}
$\times 17.5 = \qty{20.7}{MJ}$。
\textbf{9.} $72 + 20.7 = \qty{93}{MJ}$。
\textbf{10.} $72/93 = \qty{77}{\%}$（若在更完整的收支中计入\omterm{def:b1:respiration-fermentation:fermentation}{发酵}热
和结肠的能量贡献，则约为 \qty{60}{\%}；
按此计算则为四分之三）。
\textbf{11.} \qty{93}{MJ} 少于 \qty{110}{}：牛必须多吃
（\qty{12}{kg}），或吃更好的牧草，或动用\omterm{def:b1:lipids:triglyceride}{脂肪}储备。
\textbf{12.} 丙酸盐；肝脏中的\omterm{def:b1:biosyntheses-integration:gluconeogenesis}{糖异生}。
\textbf{13.} $0.9\times 2000\times 17.5 = \qty{31.5}{MJ}$。
\textbf{14.} $0.7\times 1500\times 17.5 = \qty{18.4}{MJ}$。
\textbf{15.} $0.45\times 5 = \qty{2.25}{kg}$；挥发性\omterm{def:b1:lipids:fattyacid}{脂肪酸} $0.75\times 2250\times
17.5 = \qty{29.5}{MJ}$。
\textbf{16.} $0.15\times 2250 = \qty{340}{g}$，随粪便丢失（后肠
位于小肠之后）。
\textbf{17.} $31.5 + 18.4 + 29.5 = \qty{79}{MJ}$；挥发性\omterm{def:b1:lipids:fattyacid}{脂肪酸}占比 $29.5/79 =
\qty{37}{\%}$。
\textbf{18.} 不能：它需要 $110/79\times 10 = \qty{14}{kg}$ 牧草，
只能靠每天进食十六小时，并使食物
通过速度加快一倍来维持。
\textbf{19.} 牛 $93/10 = \qty{9.3}{MJ/kg}$；马 $79/10 =
\qty{7.9}{MJ/kg}$。
\textbf{20.} 总摄入量 $10\,000\times 0.85\times 17.5 =
\qty{149}{MJ}$（不计木质素和灰分）；甲烷 \qty{7.7}{MJ} 占其
\qty{5}{\%}；一百头牛：$100\times 140\times 365 =
\qty{5.1}{t}$ 甲烷，每年。
\textbf{21.} 牛：$10\times 70/24 = \qty{29}{kg}$ 干物质存于
\omterm{prop:b1:digestion-absorption:tract}{消化道}中（另加其十倍的水分）；马：$10\times 36/24 =
\qty{15}{kg}$。
\textbf{22.} 原则上，全部 \qty{340}{g} 均可利用，价值 \qty{6}{MJ}，
并含有其必需\omterm{def:b1:proteins:aminoacid}{氨基酸}；马的盲肠内容物不像兔子那样
分离出柔软、富含营养的部分，而体重半吨的动物
无法从地面上回收自己的粪便，
使其保持可用状态。
\textbf{23.} 瘤胃微生物以尿素中的氮和秸秆中的碳合成\omterm{def:b1:proteins:peptide}{蛋白质}，
牛\omterm{def:b1:digestion-absorption:nutrition}{消化}这些微生物，几乎不需要
膳食\omterm{def:b1:proteins:peptide}{蛋白质}。马的微生物位于小肠之后，
其\omterm{def:b1:proteins:peptide}{蛋白质}随之损失，马必须
直接从食物中获取\omterm{def:b1:proteins:aminoacid}{氨基酸}。
\textbf{24.} 马\omterm{def:b1:digestion-absorption:nutrition}{吸收}葡萄糖，并将\omterm{def:b1:carbohydrates:polysaccharide}{糖原}储存于其
肌肉中，以供无氧冲刺；牛不\omterm{def:b1:digestion-absorption:nutrition}{吸收}葡萄糖，
而是从丙酸盐缓慢生成葡萄糖，并携带一个一百升的\omterm{def:b1:respiration-fermentation:fermentation}{发酵}瓮——
为放牧而非奔跑而生。
\textbf{25.} 牛的可代谢能约有四分之三
（按此收支，完整计算时为 \qtyrange{60}{70}{\%}）以挥发性\omterm{def:b1:lipids:fattyacid}{脂肪酸}形式存在，
而马则为三分之一；每千克牧草，牛为 \qty{9.3}{MJ}，
马为 \qty{7.9}{MJ}。
\end{solution}
